\chapter{Análisis melódico detallado}\label{ap:melodico}

\section{Figuras melódicas}

La voz aguda se corta en frases (notas separadas por menos de un pulso) y en cada frase se miden los intervalos entre
notas seguidas, los giros (la figura de cada nota con la anterior y la siguiente), el espectro de grados sobre el centro y
las notas extrañas al acorde del compás. Se llama grado conjunto a un intervalo de un cuarto de tono a un tono, y
microtonal a todo intervalo que no es un número entero de semitonos. La tabla~\ref{tab:figuras} los cuenta por
canción.

\begin{table}[ht]
  \centering\small
  \caption{Figuras melódicas de la voz aguda.}\label{tab:figuras}
  \begin{tabular}{@{}lrrrrrr@{}}
    \toprule
    Canción & \shortstack{Grado\\conjunto} & Saltos & Repetidas & \shortstack{Int.\\microtonales} & \shortstack{Espectro\\en \cuarto{}} & \shortstack{Notas\\extrañas} \\
    \midrule
    Sarniezz & 46 & 33 & 21 & 48 & 25 & 52 \\
    Mata Zyklek & 56 & 14 & 30 & 45 & 24 & 43 \\
    Fabienk & 43 & 33 & 24 & 23 & 14 & 39 \\
    Utzp & 42 & 24 & 34 & 25 & 12 & 37 \\
    Yor Zarad & 48 & 27 & 25 & 29 & 17 & 50 \\
    \bottomrule
  \end{tabular}
\end{table}

Aunque las notas sueltas coinciden con la partitura solo en el 46--67\,\% de los casos, los repartos coinciden: en
Sarniezz los intervalos microtonales son el 50\,\% en la herramienta y el 51\,\% en la partitura; en Mata Zyklek, 48 y
51\,\%; en Fabienk, 20 y 22\,\%. Los cuartos de tono son parte del lenguaje melódico: la partitura también tiene la mitad
de los intervalos microtonales en Sarniezz y Mata Zyklek. La nota extraña más común es la inflexión de cuarto de tono de
una nota del acorde; en Mata Zyklek y Utzp, la bordadura.

\section{Patrones de ciclo}

Para cada ciclo de cada voz, las vueltas se superponen en una rejilla de semicorcheas y el patrón es la nota que aparece
en cada casilla en al menos la mitad de las vueltas. Cada vuelta se compara con el patrón y se clasifica su variación
(inflexión de cuarto de tono, cambio de octava, notas omitidas o agregadas, transposición).

En el bajo de Mata Zyklek lo que más varía entre vueltas es Sol\sh+1 $\leftrightarrow$ Sol1, la fundamental y su segunda
neutra, tres cuartos de tono más arriba (\fl2+); en el de Yor Zarad sobre \fl II+5, el Sol resuelve hacia arriba a Sol\sh+: el grado de cuarto de
tono es la meta. En la partitura las vueltas tampoco se repiten iguales (se acogen al patrón el 60--67\,\% de las
notas): la variación entre vueltas es real, no solo error de transcripción.

\section{Centros por tramo}

La canción se corta en tramos de al menos cuatro compases, y en cada uno se mide el centro con la misma regla que el de la
canción (tabla~\ref{tab:tramos}).

\begin{table}[ht]
  \centering\small
  \caption{Centros por tramo.}\label{tab:tramos}
  \begin{tabularx}{\textwidth}{@{}lc>{\raggedright\arraybackslash}X@{}}
    \toprule
    Canción & Centro & Lo que muestran los tramos \\
    \midrule
    Sarniezz & Mi & Los tramos B giran sobre Re (\fl7) y el C sobre Sol (\fl3); Mi solo ocupa el 6\,\% de los tramos \\
    Mata Zyklek & Sol & La introducción está centrada en Sol\sh+ (\fl2+); el tramo C va a Re y vuelve a Sol por quinta \\
    Fabienk & Mi & Sin centro claro por tramo \\
    Utzp & Mi & La introducción está sobre Si (5) y entra a Mi por quinta abajo \\
    Yor Zarad & Sol & Desde 4:56 hasta el final el centro está en Sol\sh+ (\fl2+) \\
    \bottomrule
  \end{tabularx}
\end{table}

Contra la partitura, el centro por tramo coincide en el 89\,\% de los compases de Sarniezz y en el 87\,\% de los de Mata
Zyklek. A gran escala, la armonía resuelve como una cadencia de dominante a tónica entre tramos (Utzp, Mata Zyklek) o se
desplaza a un centro de cuarto de tono (Mata Zyklek, Yor Zarad), fuera del círculo de quintas.

\chapter{Arbitraje a ciegas: detalle}\label{ap:arbitraje}

Las tablas~\ref{tab:arbitraje1} y~\ref{tab:arbitraje2} detallan las dos rondas del arbitraje a ciegas
(sección~\ref{sec:evaluacion}): qué se juzgó, cuántos clips se escucharon y qué se decidió.

\begin{table}[ht]
  \centering\small
  \caption{Primera ronda del arbitraje a ciegas.}\label{tab:arbitraje1}
  \begin{tabularx}{\textwidth}{@{}l>{\raggedright}Xc>{\raggedright\arraybackslash}X@{}}
    \toprule
    Canción & Nota de la partitura & Clips & Resultado \\
    \midrule
    Sarniezz & Frase 1, posición 30: Fa\sh & 2 & Fa\sh+ en los dos (alta): cambia \\
    Sarniezz & Frase 1, posición 42: Do\sh+ & 3 & La partitura en los tres: se queda \\
    Sarniezz & Frase 1, posición 13: Fa\sh & 1 & Fa\sh+ con seguridad baja: sin decidir \\
    Sarniezz & Frase 1, posición 25: La+ & 2 & Uno la partitura y otro La\sh: sin decidir \\
    Sarniezz & Frase 1, posición 28: La+ & 1 & La\sh{} en un solo clip: sin decidir \\
    Fabienk & Patrón W de la guitarra 2, posición 3: Re\sh & 3 & Re\sh+ en los tres: cambia \\
    Mata Zyklek & Fa de la melodía 1 & 3 & La partitura en los tres: se queda \\
    Mata Zyklek & Re de la melodía 1 & 3 & La partitura en los tres: se queda \\
    \bottomrule
  \end{tabularx}
\end{table}

\begin{table}[ht]
  \centering\small
  \caption{Segunda ronda del arbitraje a ciegas.}\label{tab:arbitraje2}
  \begin{tabularx}{\textwidth}{@{}l>{\raggedright}Xc>{\raggedright\arraybackslash}X@{}}
    \toprule
    Canción & Qué se juzgó & Clips & Resultado \\
    \midrule
    Sarniezz & Frase 1, posición 13: Fa\sh & 3 & Fa\sh+ en 2 (alta): cambia \\
    Sarniezz & Frase 1, posición 25: La+ & 3 & La\sh{} en 2 (alta): cambia \\
    Sarniezz & Frase 1, posición 28: La+ & 3 & La\sh{} en 3 (alta): cambia \\
    Sarniezz & Acorde largo de la posición 19 & 2 & La partitura en los dos: se queda \\
    Sarniezz & Fase de la frase 2 de la guitarra 2 (c.~43--51) & 2 & La partitura en los dos: se queda \\
    Fabienk & Notas de la guitarra 1 en «Corrido otra vez» & 2 & «Sin la guitarra 1» en los dos (alta): se quitan \\
    Fabienk & Giro al unísono (Z, posición 3 de la guitarra 1) & 2 & «Las dos en Re\sh» en los dos (media) \\
    Mata Zyklek & Re antes del silencio, c.~120 y c.~64 & 1 + 1 & Ninguna versión: no se añade \\
    \bottomrule
  \end{tabularx}
\end{table}

Con los cambios de las dos rondas, la melodía principal pasó de 66.5 a 66.4 en Mata Zyklek, de 43.9 a 44.9 en Sarniezz y
de 45.4 a 46.2 en Fabienk; el bajo no cambió. En el c.~64 de Mata Zyklek el autor oyó un Fa agudo donde la partitura
tenía silencio; la altura medida en la pista de guitarra lo confirmó y se añadió a la guitarra 3.
